\chapter{RELATED THEORIES AND LITERATURE REVIEW}
\label{chap:litreview}

This chapter surveys the theoretical foundations and prior research directly relevant
to the design of the Vision-Language-Action (VLA) architecture proposed in this thesis.
The review is organised into five principal areas: the evolution of robotic manipulation
control paradigms, the architecture and capabilities of Vision-Language Models, the
theoretical basis of diffusion and flow-matching action generation, approaches to action
tokenisation, and a critical comparative analysis of seminal VLA architectures that
directly inform the design decisions presented in Chapter~\ref{chap:methodology}.


%─────────────────────────────────────────────────────────────────────────────
\section{ROBOTIC MANIPULATION AND LEARNING-BASED CONTROL}
%─────────────────────────────────────────────────────────────────────────────

\subsection{Classical Approaches to Robot Control}

Traditional robot manipulation systems are constructed around a modular pipeline
comprising perception, state estimation, motion planning, and low-level control
subsystems, each of which is engineered independently and integrated through
well-defined interfaces. Perception modules employ geometric models and hand-designed
feature detectors to identify object poses from depth or RGB sensor data; planning
modules invoke algorithms such as Rapidly-exploring Random Trees (RRT) or
Probabilistic Roadmaps (PRM) to compute collision-free trajectories in configuration
space; and control modules execute these trajectories via Proportional-Integral-Derivative
(PID) regulators or model-predictive controllers~\cite{lavalle2006planning}.

These pipelines achieve high reliability within the specific operational envelope for
which they are engineered. However, their dependence on explicit world models, manually
specified object representations, and task-specific heuristics renders them brittle in
the face of distributional change: the introduction of a previously unseen object
geometry, an alteration in lighting conditions, or a novel task goal typically
necessitates substantial re-engineering. This fundamental limitation motivates the
shift toward data-driven, learning-based approaches.

\subsection{Imitation Learning and Behaviour Cloning}

Imitation learning addresses the robot control problem by learning a policy
$\pi_\theta(a \mid o)$ that maps observations $o$ to actions $a$ by supervised
regression against a dataset of expert demonstrations $\mathcal{D} = \{(o_i, a_i)\}$.
In its simplest instantiation, known as Behaviour Cloning (BC), the policy is trained
by minimising the mean squared error between predicted and demonstrated
actions~\cite{pomerleau1989alvinn}. Despite its conceptual simplicity, BC is susceptible
to compounding errors arising from covariate shift: small deviations from the
demonstration distribution at inference time produce states that are increasingly
out-of-distribution, leading to cascading failures. Methods such as DAgger~\cite{ross2011dagger}
address this by iteratively augmenting the training distribution with on-policy
observations relabelled by the expert, though they impose a requirement for interactive
expert access that is often prohibitive in physical robot settings.

\subsection{End-to-End Visuomotor Policies}

The introduction of deep convolutional neural networks as visual encoders enabled
the development of end-to-end visuomotor policies that directly map raw pixel
observations to motor commands without explicit intermediate representations.
Seminal work by Levine et al.~\cite{levine2016end} demonstrated that such policies
could be trained from large collections of robot demonstrations to perform a variety
of manipulation tasks, albeit within a restricted range of objects and scene
configurations. Subsequent work established that architectural choices, particularly
the design of the visual encoder, the action output head, and the temporal context
window, significantly influence the generalisation capacity of the resulting
policy. This body of work provides the direct precursor to the VLA paradigm, in
which the visual encoder and policy backbone are replaced by a pre-trained foundation
model of substantially greater representational capacity.


%─────────────────────────────────────────────────────────────────────────────
\section{VISION-LANGUAGE MODELS}
%─────────────────────────────────────────────────────────────────────────────

\subsection{Transformer Architecture}

The Transformer~\cite{vaswani2017attention} is the foundational architectural
component underlying all modern VLMs. It processes input sequences through a
series of multi-head self-attention layers, each of which computes a weighted
aggregation of value vectors for every position in the sequence, with attention
weights determined by the scaled dot-product between query and key vectors:

\begin{equation}
    \text{Attention}(Q, K, V) = \text{softmax}\!\left(\frac{QK^\top}{\sqrt{d_k}}\right)V
    \label{eq:attention}
\end{equation}

where $Q$, $K$, and $V$ are linear projections of the input, and $d_k$ is the
dimension of the key vectors. The self-attention mechanism enables each position
to attend to all other positions in the sequence, endowing the model with the
capacity to capture long-range dependencies that recurrent architectures struggle
to represent efficiently. Position-wise feed-forward sublayers, residual connections,
and layer normalisation are interleaved with the attention sublayers to form a
complete Transformer block. Large language models are constructed by stacking
many such blocks and pre-training on internet-scale text corpora via the next-token
prediction objective.

\subsection{Vision Encoders: SigLIP and DINOv2}

Integrating visual inputs into language models requires a vision encoder that maps
raw image pixels to a sequence of embedding vectors compatible with the language
backbone's representation space. Two families of vision encoders are particularly
relevant to the present work.

\textbf{SigLIP}~\cite{zhai2023sigmoid} is a contrastive vision-language encoder
trained with a sigmoid loss function over image-text pairs. Unlike the softmax-based
contrastive objective of CLIP~\cite{radford2021clip}, which requires computing
pairwise similarities across all examples in a batch, the sigmoid formulation treats
each image-text pair independently as a binary classification problem. This confers
both computational advantages at large batch sizes and improved performance on
tasks requiring fine-grained semantic alignment between visual and linguistic
representations. SigLIP is the vision encoder employed in the VLM backbone of the
proposed architecture.

\textbf{DINOv2}~\cite{oquab2023dinov2} is a self-supervised vision encoder trained
by distilling features from a teacher network using a combination of image-level and
patch-level objectives. Because DINOv2 is trained without image-text supervision,
its feature representations are less constrained by semantic label boundaries and
exhibit notably strong spatial structure: patch-level features encode precise
geometric and texture information at the scale of individual image regions. This
property makes DINOv2 well-suited as a dedicated end-effector vision encoder, where
fine-grained spatial features of the immediate manipulation workspace are more
informative than high-level semantic labels. The proposed architecture employs a
successor model, DINOv3, in this role.

\subsection{Multimodal Language Models and Instruction Following}

Modern VLMs such as LLaVA~\cite{liu2023llava} and GPT-4V~\cite{openai2023gpt4v}
extend the language Transformer backbone to accept interleaved sequences of image
tokens and text tokens. Visual inputs are first encoded by a vision encoder into a
sequence of patch embeddings, which are then projected into the language model's
token embedding space via a learned linear or MLP adapter. The resulting multimodal
sequence is processed by the language backbone in the same manner as a purely
textual sequence, enabling the model to reason jointly about visual and linguistic
content. Pre-training on large-scale image-text corpora, followed by instruction-tuning
on curated demonstrations of instruction-following behaviour, yields models that can
respond to complex compositional queries about image content with remarkable
accuracy and generalisation.

The instruction-following capability of these models is directly relevant to
robot manipulation: if a VLM can correctly parse an instruction such as
\textit{``put the banana in the green box''} and localise the referenced objects
within a scene image, it provides a strong semantic grounding prior upon which
an action generation module can be conditioned.

\subsection{Object Localisation as Spatial Grounding}

A dedicated line of research has investigated the extent to which VLMs can be
trained to produce spatially grounded outputs in the form of explicit bounding box
coordinates~\cite{chen2024localization}. In this paradigm, the model is trained
to generate token sequences such as
\texttt{banana <loc12><loc25><loc67><loc89>}, where the
\texttt{<loc>} tokens index a discretised coordinate grid. Training on this
localisation objective has been demonstrated to substantially improve the spatial
reasoning capacity of the underlying language backbone, with downstream benefits
for tasks that require precise object-relative positioning---a property of direct
relevance to manipulation planning. Accordingly, object localisation is incorporated
as an auxiliary pre-training task in the proposed architecture.


%─────────────────────────────────────────────────────────────────────────────
\section{DIFFUSION MODELS AND FLOW MATCHING}
%─────────────────────────────────────────────────────────────────────────────

\subsection{Denoising Diffusion Probabilistic Models}

Denoising Diffusion Probabilistic Models (DDPMs)~\cite{ho2020ddpm} define a
generative model through a forward process that progressively corrupts a data
sample $x_0$ with Gaussian noise over $T$ timesteps:

\begin{equation}
    q(x_t \mid x_{t-1}) = \mathcal{N}\!\left(x_t;\, \sqrt{1 - \beta_t}\, x_{t-1},\, \beta_t I\right)
    \label{eq:ddpm_forward}
\end{equation}

and a learned reverse process that denoises the corrupted sample step by step:

\begin{equation}
    p_\theta(x_{t-1} \mid x_t) = \mathcal{N}\!\left(x_{t-1};\, \mu_\theta(x_t, t),\, \Sigma_\theta(x_t, t)\right)
    \label{eq:ddpm_reverse}
\end{equation}

The reverse process is parameterised by a neural network trained to predict
the noise $\epsilon$ added at each step, minimising the objective:

\begin{equation}
    \mathcal{L}_\text{DDPM} = \mathbb{E}_{x_0, \epsilon, t}\!\left[\left\| \epsilon - \epsilon_\theta(x_t, t) \right\|^2\right]
    \label{eq:ddpm_loss}
\end{equation}

DDPMs have achieved state-of-the-art results in image generation and have been
subsequently applied to robot action generation, where the data distribution to
be modelled is the distribution of expert trajectory chunks rather than natural
images~\cite{chi2023diffusionpolicy}.

\subsection{Flow Matching}

Flow Matching~\cite{lipman2022flow} is a more recent and computationally efficient
alternative to DDPM-based generation. Rather than defining a multi-step noising
process and learning to reverse it, flow matching directly learns a vector field
$v_\theta(x, t)$ that transports samples from a simple source distribution
(typically $\mathcal{N}(0, I)$) to the target data distribution along straight-line
paths in probability space:

\begin{equation}
    \frac{dx}{dt} = v_\theta(x, t), \quad x(0) \sim \mathcal{N}(0, I), \quad x(1) \sim p_\text{data}
    \label{eq:flow}
\end{equation}

The vector field is trained by regressing against the conditional flow between
pairs of noise and data samples. Compared to DDPM, flow matching admits fewer
function evaluations at inference time, leading to substantially lower action
generation latency---a critical property for real-time robot control. The proposed
architecture employs a flow-matching objective within its Diffusion Transformer
action expert.

\subsection{Diffusion Transformer (DiT)}

The Diffusion Transformer (DiT)~\cite{peebles2023scalable} replaces the U-Net
backbone traditionally employed in diffusion models with a Transformer architecture,
enabling the model to condition on arbitrary context vectors via cross-attention
or adaptive layer normalisation (adaLN). This architectural choice is particularly
advantageous in the VLA setting, where the action expert must be conditioned on
heterogeneous context signals of variable length---including VLM hidden state
sequences and dense visual feature maps from the end-effector encoder. The DiT's
attention mechanism provides a principled and flexible mechanism for attending
to these conditioning signals during the denoising process.


%─────────────────────────────────────────────────────────────────────────────
\section{ACTION TOKENISATION}
%─────────────────────────────────────────────────────────────────────────────

\subsection{Discrete Action Representations}

A central challenge in adapting language model backbones to robot control is the
mismatch between the discrete token vocabulary of language models and the continuous,
high-dimensional nature of robot joint trajectories. One family of approaches
resolves this mismatch by discretising continuous trajectories into a finite
vocabulary of action tokens, enabling the language backbone to predict actions
autoregressively using the same next-token prediction objective employed during
language pre-training. Early work in this direction employed na\"{i}ve uniform
binning of individual joint angles~\cite{brohan2022rt1}, but this approach incurs
high quantisation error and an exponentially large vocabulary for multi-joint
trajectories.

\subsection{FAST+ Tokeniser}

The Frequency-space Action Sequence Tokeniser (FAST)~\cite{pertsch2025fast}
addresses these limitations by representing trajectory chunks in the frequency
domain via a Discrete Cosine Transform (DCT) prior to quantisation. The key
insight is that robot trajectories are temporally smooth signals whose energy
is concentrated in low-frequency DCT coefficients; retaining only the dominant
coefficients before quantisation therefore preserves the essential kinematic
content of the trajectory with substantially fewer tokens and less quantisation
error than spatial-domain methods. The FAST+ variant further improves upon
the original FAST tokeniser through refinements to the codebook construction
and decoding procedures, and has been empirically validated as the
state-of-the-art discrete action tokeniser across a diverse set of
manipulation benchmarks. Accordingly, FAST+ is adopted as the primary action
tokenisation strategy in the proposed architecture, in accordance with current
best practice.

\subsection{Continuous State Representations}

An alternative to discrete tokenisation is to represent robot state and action
information as continuous vectors and inject them into the action model through
learned conditioning mechanisms. Feature-wise Linear Modulation
(FiLM)~\cite{perez2018film} is one such mechanism: it conditions a neural network
layer on an external vector by learning affine transformation parameters
$(\gamma, \beta)$ that modulate the layer's feature activations:

\begin{equation}
    \text{FiLM}(h \mid z) = \gamma(z) \odot h + \beta(z)
    \label{eq:film}
\end{equation}

where $h$ is the layer's activation and $z$ is the conditioning vector.
FiLM injection is employed by the GR00T N1 architecture~\cite{nvidia2024groot}
to condition the action expert on proprioceptive state observations, and is noted
in the proposed architecture as an alternative conditioning pathway to hidden state
injection. The relative merits of these two conditioning strategies are empirically
evaluated in Chapter~\ref{chap:experiments}.


%─────────────────────────────────────────────────────────────────────────────
\section{RELATED VISION-LANGUAGE-ACTION ARCHITECTURES}
%─────────────────────────────────────────────────────────────────────────────

\subsection{OpenVLA}

OpenVLA~\cite{kim2024openvla} is an open-source VLA model that established
several widely adopted architectural conventions. It employs a dual-encoder
vision backbone combining SigLIP for semantic visual features and DINOv2 for
spatially structured patch features, whose outputs are concatenated and projected
into the token embedding space of a pre-trained language model. Action outputs
are produced by appending action token predictions to the instruction-grounded
visual context, trained end-to-end via next-token prediction. OpenVLA demonstrated
that fine-tuning a general-purpose VLM backbone on robot demonstration data
yields a manipulation policy with substantially greater generalisation to novel
objects and instructions than task-specific baselines. However, its reliance on
a purely discrete autoregressive action output without a dedicated continuous
action refinement module limits its trajectory smoothness and precision on tasks
requiring fine-grained motor control.

\subsection{GR00T N1}

GR00T N1~\cite{nvidia2024groot} is NVIDIA's open foundation model for humanoid
robot control. It adopts a heterogeneous architecture in which a VLM backbone
processes visual and linguistic inputs, and a separate action expert is conditioned
on continuous state representations injected via FiLM. In contrast to the best
practice of employing FAST+ for action tokenisation, GR00T N1 employs a customised
tokenisation scheme optimised for the high-dimensional action spaces of full-body
humanoid robots. The architecture targets large-scale, whole-body manipulation and
locomotion tasks rather than the tabletop manipulation setting addressed in this
thesis. GR00T N1 is nonetheless a relevant baseline and point of comparison, as
it represents the most capable open VLA system currently available and embodies
a comprehensive set of architectural design choices---including the frozen backbone,
FiLM-based conditioning, and multi-camera visual input---that directly intersect
with the design space explored in the present work.

\subsection{Diffusion-VLA}

Diffusion-VLA~\cite{wen2024diffusionvla} is a VLA architecture that incorporates
a self-generated chain-of-thought reasoning module as an interpretable intermediary
between language understanding and action generation. Before producing an action,
the model generates a natural language rationale describing its interpretation of
the scene and the sub-goals required to accomplish the instruction. This rationale
is then provided as additional context to the diffusion-based action expert, enabling
more interpretable and compositionally generalisable policies. Diffusion-VLA
demonstrates that the language backbone of a VLA system can be leveraged not only
as a feature extractor but as an active reasoning module, though the additional
inference latency incurred by rationale generation is a practical limitation for
time-sensitive control loops.

\subsection{Knowledge Insulating VLA}

The Knowledge Insulating VLA framework~\cite{pertsch2025kilovla} proposes a
pre-training strategy specifically designed to prevent the catastrophic forgetting
of general visual and linguistic knowledge during robot-specific fine-tuning.
The central observation is that na\"{i}ve fine-tuning of a VLM backbone on
robot demonstration data degrades the backbone's generalisation capacity,
because gradient updates from the manipulation task overwrite the broadly useful
representations acquired during foundation model pre-training. The proposed
remedy---selectively freezing the backbone and training only lightweight adapter
modules alongside the action expert---is demonstrated to accelerate convergence,
reduce data requirements, and improve generalisation to novel objects and
instructions relative to full fine-tuning baselines. This finding motivates
the frozen-backbone design adopted in the present architecture.

\subsection{Summary and Research Gap}

Table~\ref{tab:vla_comparison} summarises the key architectural and training
characteristics of the surveyed VLA models in relation to the proposed architecture.

\begin{table}[ht]
\centering
\caption{Comparison of related VLA architectures. FM = Flow Matching,
         FS = Frozen Backbone, CT = Continuous Tokenisation,
         DT = Discrete Tokenisation.}
\label{tab:vla_comparison}
\begin{tabular}{lcccccc}
\hline
\textbf{Model} & \textbf{Vision Enc.} & \textbf{Action Gen.} & \textbf{Tokeniser}
               & \textbf{FS} & \textbf{End-Eff. Cam.} & \textbf{Localisation} \\
\hline
OpenVLA
  & SigLIP + DINOv2 & Autoregressive & Uniform bin
  & \xmark & \xmark & \xmark \\
GR00T N1
  & SigLIP          & Diffusion (FM) & Custom (CT)
  & \cmark & \cmark & \xmark \\
Diffusion-VLA
  & SigLIP          & Diffusion (FM) & FAST+ (DT)
  & \cmark & \xmark & \xmark \\
KILO-VLA
  & SigLIP + DINOv2 & Autoregressive & FAST+ (DT)
  & \cmark & \xmark & \xmark \\
\textbf{Ours}
  & SigLIP + DINOv3 & Diffusion (FM) & FAST+ (DT)
  & \cmark & \cmark & \cmark \\
\hline
\end{tabular}
\end{table}


The comparison reveals a consistent gap across all surveyed architectures: no
existing model simultaneously employs a dedicated end-effector vision expert,
FAST+ discrete action tokenisation, a flow-matching diffusion action expert, and
object localisation as an auxiliary pre-training task within a frozen-backbone
framework. The proposed architecture is designed to occupy this gap, synthesising
the respective strengths of the surveyed systems into a unified design.